\documentclass[11pt, letterpaper]{article}
\usepackage[spanish]{babel}
\usepackage[utf8]{inputenc}
\usepackage{amsmath,amssymb,amsfonts}
\usepackage{enumerate}
\usepackage{subfig}
\usepackage{nicefrac}
\usepackage{url}
\usepackage{multirow}
\usepackage{graphicx}
\usepackage{setspace}

\usepackage[right=2.5cm, left=2.5cm, top=2cm, bottom=2cm]{geometry}

\decimalpoint

\begin{document}
% \maketitle

\thispagestyle{empty}

\noindent 
\begin{center}
  % \includegraphics[width=5em]{cebolla.pdf}\\
  Universidad Simón Bolívar\\
  Decanato de Estudios de Postgrado\\
  Postgrado en Ciencias de la Computación\\
  Representación del Conocimiento \\
  \vfill
  {\Huge Pensar un titulo}
  \vfill
  Daniel Izquierdo \\
  Julio Castillo\\[2em]
  Abril de 2009
\end{center}

\newpage

\section{Introducción}
Este documento presenta una breve descripción de cómo calcular el MPE
de una red bayesiana grande utilizando como heurística el MPE de una
red simplificada que proporciona una cota superior. El cálculo del MPE
en la red simplificada se calcula utilizando una codificación en un
circuito aritmético que se obtiene a partir de una compilación a
d-DNNF.

% Por último se muestra los resultados de una serie de experimentos que
% se realizaron para probar 

\section{Implementación}

\subsection{Cálculo de MPE utilizando teorías en d-DNNF}
\label{sec:condificacion}
Hablar aqui un poquito de las funciones multilineales

La red bayesiana se convirtió a una teoría $\Delta$ en CNF de la
siguiente forma
\begin{itemize}
\item Para cada variable $X$ de la red bayesiana que toma valores
  $x^1,\dots,x^n$, se agregan a $\Delta$ las siguiente clausulas
  \begin{align}
    \lambda_{x^1} \lor \cdots \lor \lambda_{x^n} \label{eq:cnf1} \\
    \neg\lambda_{x^i} \lor \neg\lambda_{x^j} \quad i\neq j \label{eq:cnf2}
  \end{align}
\item Para cada variable $X$ con padres $U$, se crea una variable
  $\theta_{x|u_1,\dots,u_m}$, $\Delta$
  contiene:
  \begin{align}
    \lambda_{x} \land \lambda_{u_1} \land \cdots \land \lambda_{u_m} \Rightarrow \theta_{x|u_1,\dots,u_m} \label{eq::cnf3}\\
    \theta_{x|u_1,\dots,u_m} \Rightarrow \lambda_{x}, \;\; \theta_{x|u_1,\dots,u_m} \Rightarrow \lambda_{u_1} ,\;\; \dots \;\;, \theta_{x|u_1,\dots,u_m} \Rightarrow \lambda_{u_m} \label{eq:cnf4}
  \end{align}
\end{itemize}

Posteriormente la $\Delta$ se compila a una teoría $\Gamma$
equivalente en d-DNNF smooth utilizando
\texttt{c2d}~{\cite{darwicheAAAI02}}. Luego en $\Gamma$ se sustituyen
las los literales negativos por $1$, las conjunciones por
multiplicaciones y las disjunciones por sumas. Esta sustitución
permite hacer inferencia sobre la red bayesiana sustituyendo en
$\Gamma$ las variables $\theta$ por las probabilidades
correspondientes y las variables $\lambda_x$ por 1 si $x$ es
compatible con la evidencia o 0 si es imcompatible~\cite{Darwiche01alogical}.

Esta sustitución produce lo que se conoce como una representación en
circuito aritmético de la red bayesiana. Este circuito aritmético, al
evaluarlo de la forma indicada, corresponde a evaluar la función
multilineal asociado a la red bayesiana.

% Por ejemplo, para una red con variables $A$, $B$ y $C$, la función
% multilineal asociada es
% \begin{align}
%   f & = \lambda_a\lambda_b\lambda_c\
% \end{align}

Cada sumando en la función multilineal de una red corresponde a la
probabilidad de una instanciación particular compatible con la
evidencia. Si sustituimos las sumas por el operador $\max$,
obtendremos la asignación de variables de máxima para una evidencia,
es decir, la probabilidad del MPE. Por lo tanto, si en $\Gamma$
sustituimos las disjunciones por $\max$, entonces podemos calcular el
MPE de la red utilizando el circuito aritmético.

\subsubsection{Restricciones Lógicas}
Una restricción lógica es una probabilidad condicional que es igual a
0 o 1. En cada uno de estos casos se pueden hacer simplificaciones
sobre $\Delta$ que permiten reducir significativamente el tamaño de la
teoría compilada en d-DNNF.

Si tenemos una red bayesiana con parámetro $\theta_{x|u_1,\dots,u_m}$
que toma valor cero, entonces las clausulas \eqref{eq::cnf3} y
\eqref{eq:cnf4} que involucran a $\theta_{x|u_1,\dots,u_m}$ se pueden sustituir por
\begin{equation}
  \label{eq:simpcero}
  \neg \lambda_{x} \lor \neg \lambda_{u_1} \lor \cdots \lor \neg \lambda_{u_m}
\end{equation}
Es decir, que $\theta_{x|u_1,\dots,u_m}$ se puede eliminar de
$\Delta$. Por otro lado, cuando $\theta_{x|u_1,\dots,u_m}=1$, podemos
omitir por completo a $\theta_{x|u_1,\dots,u_m}$ y las clausulas que
la mencionan.

Como se mostrará más adelante, estas simplificaciones tienen un gran
impacto sobre el tamaño de la teoría comipliada cuando la red bayesiana
esta compuesta por una gran cantidad de nodos determinísticos.

Existe un segundo tipo de simplificación conocida como
\textsl{context-specific independence}, sin embargo esta
simplificación no se implemento para este proyecto.

% \subsubsection{Detalles de implementación}
% \label{sec:detall-de-impl}
% Para hacer la traducción a CNF primero se hace un mapeo 

\subsection{Cálculo de MPE para redes grandes}
\label{sec:calculo-de-mpe}
Para redes bayesianas con muchos nodos, el paso de compilación a
d-DNNF que se explicó anteriormente no terminará en una tiempo
razonable. En este caso utilizamos la técnica de \textsl{node
  splitting} propuesta en~\cite{ChoiChaviraDarwiche07}.

La idea es tomar la red original y hacer una serie splittings hasta
que sea posible compilar la red a d-DNNF y luego utilizar el MPE sobre
esta nueva red (que se puede calcular utilizando el circuito
aritmético) como heurística para obtener el MPE de la red original.

Este método se basa en el teorema~1 de \cite{ChoiChaviraDarwiche07} el
cual dice que si $N$ es un red bayesiana y $N'$ es el resultado de
hacer un splitting en $N$, entonces
\begin{equation}
  \label{eq:2}
  MPE_p(N,\mathbf{e}) \leq \beta MPE_p(N', \mathbf{e}, \vec{\mathbf{e}})
\end{equation}
donde $\vec{\mathbf{e}}$ es el conjunto de instanciaciones compatibles con $e$,
y $\beta = \prod_{C \in \mathbf{C}} |C|$ donde $\mathbf{C}$ es el conjunto de
clones en $N'$.

Si tenemos una red bayesiana $N'$ que es el resultado de hacer
splitting sobre la red $N$, entonces el procedimiento para calcular el
$MPE(N,\mathbf{e})$ utilizando $N'$ se describe a continuación. Se
tiene una variable $\mathbf{z}$ que representa una asignación parcial
en $N$ y una segunda variable $q*$ que representa la cota inferior
actual para el MPE de de $N$. La búsqueda se inicia con $\mathbf{z} =
\mathbf{e}$ y $q* = 0$. En cada nodo del algoritmo de búsqueda
computamos $q$, una cota superior para el MPE de $N'$ y la
instanciación parcial $\mathbf{z}$, esto se puede hacer evaluando el
circuito aritmético. Si $q>q*$ entonces debemos continuar con la
búsqueda porque puede ser posible que $z$ se pueda extender para
conseguir una mejor aproximación que la que ya tenemos. Si $z$ ya
tiene todas la variables instaciadas, entonces se puede demostrar que
$Pr(z) = q$ y entonces tenemos una nueva cota. Si $z$ no es una
instanciación completa, seleccionamos una variable $X$ que no esté
instanciada y para cada valor $x$ que $X$ puede tomar, agregamos a $z$
la asignación $\{ X = x \}$ y llamamos al procedimiento recursivamente.

PONER AQUI EL ALGORITMO DE B\&B

Llamemos a $L$ el conjunto de variables de $N$ a las que se le hizo
split. Un resultado importante de~\cite{ChoiChaviraDarwiche07} es que
si todas las variables de $L$ están instanciadas entonces $MPE_p(N,
\mathbf{e}) = MPE_p(N', \mathbf{e}, \vec{\mathbf{e}})$. Como veremos
más adelante, ordenar los nodos para aprovechar esta propiedad permite
mejorar el proceso de búsqueda.

\section{Resultados}

\section{Conclusiones}


\bibliographystyle{plain}
\bibliography{informe}

\end{document}

%%% Local Variables: 
%%% mode: latex
%%% TeX-master: t
%%% End: 

